% Source: 50_Chapter_38_To_Earth_s_Remotest_People_Protestant_M.txt
\bookchapter[To Earth’s Remotest People]{CHAPTER 38}{TO EARTH’S REMOTEST PEOPLE}{Protestant Missions}

\leadin{I}{N AN ENGLISH VILLAGE,} late in the eighteenth century, stood a humble workshop. Over its
door a sign announced, “Secondhand shoes bought and sold.” Inside, the shoemaker, William Carey,
repaired a neighbor’s boot or, when time allowed, continued his study of Latin and Greek. Over the
workbench was a crude map of the world. On it Carey had penned bits of information from the voyages
of Captain James Cook or some other world traveler. A friend, Thomas Scott, called the workshop
“Carey’s College.”

Carey’s workbench and map are fit symbols of the awakening interest in distant peoples during the
Age of Progress and in the means of getting the gospel to them. At the beginning of the nineteenth
century, Protestant Christianity scarcely existed outside Europe and America. Asia was almost
untouched by the gospel, except for small traces in India and in the East Indies where the Dutch had
taken over from the Portuguese. Africa was the “dark continent” except for the ancient Copts in
Egypt and Ethiopia. After eighteen centuries Christianity was far from being a world religion.

It is a different story today. The front page of almost any metropolitan newspaper carries news,
daily, of events vitally linked with Christianity somewhere on the earth.

The great era of Christian expansion was the nineteenth century. “Never had any other set of ideas,
religious or secular, been propagated over so wide an area by so many professional agents maintained
by the unconstrained donations of so many millions of individuals.” That is the informed judgment of
Kenneth Scott Latourette, the foremost historian of Christianity’s expansion. For sheer magnitude
the Christian mission in the nineteenth century is without parallel in human history.

How do we explain this sudden explosion of Protestant energy aimed at winning the world for Christ?

\section{THE PIONEER IN MODERN MISSIONS}

During the first century of Protestant history, the Roman Catholic countries, Spain and Portugal,
dominated the commercial and imperial expansion of European peoples. The great missionary names were
Xavier, Las Casas, and Ricci. Only after the English defeat of the Spanish Armada (1588) and the
emergence of the British and Dutch as colonial powers do new continents and peoples open to
Protestant missionaries.

The first Protestants to attempt to reach distant peoples with the gospel were the Pietists.
Moravian concern, however, was focused on individuals in some European colony perishing without the
knowledge of Christ. The Christian groups created by Pietists were tiny islands in the surrounding
sea of “heathenism.”

William Carey introduced Christians to missions on a grander scale. He thought in terms of the
evangelization of whole countries and of what happens when whole populations become Christian. He
held that the foreign missionary can never make more than a small contribution to the accomplishment
of the work that has to be done, and that therefore the development of the local ministry is the
first and greatest of all missionary considerations. Above all, he saw that Christianity must be
firmly rooted in the culture and traditions of the land in which it is planted. For all these
reasons and more Carey gained the title “Father of Modern Missions.”

The English cobbler was a most unlikely candidate for greatness. He was married to a poor,
uneducated girl, and what Carey earned as a shoemaker was often too little to provide enough to eat.
Yet the greatness of the man was within, not in his circumstances. He had a ravenous hunger for
knowledge and would go without food to buy a book. Columbus and Captain James Cook were his great
heroes.

In 1779, through a fellow shoemaker, he was converted to faith in Christ, and in 1783 he was
baptized as a believer. After gaining some preaching experience he became pastor of the Moulton
Baptist Chapel, supporting his family by teaching and shoemaking.

In Baptist circles he met fellow pastor Andrew Fuller. Fuller, a strong Calvinist, broke with some
of his fellow Calvinists who pictured vigorous evangelism and appeals to conversion as inconsistent
with God’s election of certain individuals to salvation. Fuller held to both his belief in a
Calvinist version of election and the mandate to follow Jesus and the apostles in the practice of
evangelism.

“We have sunk into such a compromising way of dealing with the unconverted,” Fuller complained, “as
to have well nigh lost the spirit of the primitive preachers, and hence it is that sinners of every
description can sit so quietly as they do, year after year, in our places of worship.”

From Fuller’s teaching Carey drew the inescapable inference that if it is the duty of all men to
repent and believe the gospel, as Fuller argued, it is also the duty of those entrusted with the
gospel to carry it to the whole world.

In 1792 Carey published An Enquiry into the Obligation of Christians to Use Means for the Conversion
of the Heathen. It created an epoch. In it Carey took up the five objections that people raised
against missions to heathen lands: their distance, their barbarism, the danger that would be
incurred, the difficulties of support, and the unintelligible languages. One by one he answered
these. The same obstacles had not prevented the merchants from going to distant shores. “It only
requires,” he wrote, “that we should have as much love to the souls of our fellow-creatures, and
fellow sinners, as they have for the profits arising from a few otter skins, and all these
difficulties could be easily surmounted.” He ended his appeal with practical proposals for the
preaching of the gospel throughout the world.

By encouraging each other Carey and Fuller succeeded in breaking free from the restrictive theology
of their time. They went back to the New Testament, especially to Jesus’ injunction to preach the
gospel to all the world and to the apostle Paul’s declaration of God’s intention: “that at the name
of Jesus every knee should bow, in heaven and on earth and under the earth, and every tongue confess
that Jesus Christ is Lord, to the glory of God the Father” (Phil. 2: 10, 11). Such texts were read
with a sense of immediacy, as if Jesus were speaking to them and not merely the disciples long ago.

As a result, in October 1792 Carey, Fuller, and eleven Baptist colleagues formed the Baptist
Missionary Society, and within a year Carey and his family were on their way to India. His wife
Dorothy had planned to stay behind, but she reversed her decision and accepted a final appeal to
join her husband. Almost immediately, she regretted her decision. Eventually she lived out her life
in her room in bitterness and perhaps insanity.

The British East India Company, which had been the virtual ruler of India since 1763, was exercising
its full power at that time. It was not enthusiastic about missions. Its interest was in profits.
Most of its representatives, living free and easy lives and enjoying to the full their sense of
racial superiority, considered “the sending out of missionaries into our Eastern possessions to be
the maddest, most extravagant, most costly, most indefensible project which has ever been suggested
by a moonstruck fanatic. Such a scheme is pernicious, imprudent, useless, harmful, dangerous,
profitless, fantastic.”

The Company refused Carey permission to live in Calcutta, so he settled instead in Serampore, under
the Danish flag. He secured employment as foreman of an indigo factory in Bengal. Since the position
demanded only three months of the year, he found plenty of time for intensive study of the oriental
languages. In 1799 two fellow Baptists, Joshua Marshman and William Ward, joined Carey at Serampore.
For the next quarter-century the three men worked together to organize a growing network of mission
stations in and beyond Bengal.

Carey and his companions plunged courageously into all the intricacies of Hindu thought. They did
not regard these studies as in any way a distraction from their missionary work. On the contrary,
they regarded a full understanding of Hindu thought as an essential part of their equipment, not
only because the preacher of the gospel cannot be clearly understood if he speaks merely out of the
self-confidence of his own knowledge, but also because they understood that it is not only the souls
and bodies of men that need to be redeemed; the thought-world of a non-Christian nation is also one
of those realms that are to be taken captive and brought into subjection to Christ. By 1824 Carey
had supervised six complete and twenty-four partial translations of the Bible, as well as publishing
several grammars, dictionaries, and translations of Eastern books.

\section{THE CONTAGION OF MISSIONARY SERVICE}

The example of the Serampore trio proved contagious. The beginning of the nineteenth century found a
new and pervasive determination in Protestantism to carry the gospel to all men. The earlier
prevailing attitude of the major churches had considered missions an unnecessary and hopeless
undertaking. Now voices were raised on all sides proclaiming the duty of all Christians to share in
the conversion of the peoples of the whole world. The gospel was not the private possession of
European peoples.

The list of missionary pioneers would run into the hundreds: Henry Martyn in India, Robert Morrison
in China, John Williams in the South Seas, Adoniram Judson in Burma, Alexander Duff in India, Allen
Gardiner on Tierra del Fuego, Robert MoVat in South Africa, and many more. Scores of other
missionaries and their wives are long forgotten because they died in a matter of months in some
malaria-infested tropical climate or at the hands of some savage tribe.

In large part this new passion to preach the gospel to the heathen sprang from those portions of
Protestantism deeply influenced by the eighteenth-century evangelical revivals in England and
America. For the first three decades of the new missionary era the endeavor was almost exclusively
evangelical.

This is hardly surprising. The evangelical awakenings revolutionized preaching and its objectives.
Traditional churchmen usually limited the minister’s task to nurturing the seed of faith planted at
baptism in virtually all members of the parish. Such men could not imagine preaching the gospel in a
tribal society. At the same time those Christians who held a rigid doctrine of predestination never
seemed to concern themselves with the elect in India or China. Evangelicals, however, like Carey,
saw preaching as calling sinners to God through faith in Christ. They felt a personal responsibility
to do this and saw no difference in principle between “baptized heathen” in Britain and
non-Christian peoples overseas.

Only in the 1820s and 1830s did interest in overseas missions become a general feature of British
church life. This was due in part to the success of Evangelicals in influencing English and Scottish
society. Many of their values were adopted outside their circle. In particular the idea of Britain
as a Christian nation, with Christian responsibilities overseas, took root.

The vision of this missionary task is expressed in thousands of sermons and hundreds of hymns of
that time. Reginald Heber, who gave his life for India serving as Bishop of Calcutta, provides an
example in his widely sung “From Greenland’s Icy Mountain”:

\begin{poem}
Can we, whose souls are lighted\\
With wisdom from on high;\\
Can we to men benighted\\
The lamp of life deny?\\
Salvation, O salvation!\\
The joyful sound proclaim,\\
Till earth’s remotest nation\\
Has learnt Messiah’s Name.
\end{poem}

Two emphases led to this new Protestant world vision. One, as Carey and Fuller illustrate, was
evangelical. The Bible teaches that men are lost without faith in Christ, and the Lord commands
believers in every age to make salvation known in all the world.

The other was prophetic. Many Christians in the nineteenth century followed Jonathan Edwards in the
belief that the knowledge of the Lord would fill the earth as the waters cover the sea, and this
spread of the gospel was preparation for the coming reign of Christ upon the earth. This belief in a
future reign of Christ was called millennialism.

The Protestant mission to all the world was no empty dream. The dedication of the missionary
movement blended with the optimism of the Age of Progress to make the achievement of the goal seem
quite possible. Thus the Student Volunteer Movement for Foreign Missions took as its watchword, “the
evangelization of the world in this generation.”

The vision was constantly renewed by some fresh account from Africa or the South Pacific. None of
these proved more inspiring than the reports of the spiritual darkness of Africa or the horrors of
the Arab slave trade sent home by David Livingstone (1813–1873).

The great explorer of the dark continent came from a hardy clan of Scotsmen. When he was nineteen,
he determined to devote his life to the “alleviation of human misery.” He studied as a doctor to
prepare himself for the work of a missionary, and, attracted by the fame of Robert MoVat in South
Africa, he went to Africa to help in the work.

\section{LIVINGSTONE IN AFRICA}

Arriving in 1841, Livingstone served for ten years in the ordinary routine of missionary work. But
he was not a man to stay long in any one place. The mind and impulse of the explorer were in him,
and he was always drawn on, in his own words, by “the smoke of a thousand villages” that had never
seen a missionary.

The first great journey that made him famous led him through the jungles to the west coast in
Angola, and then, because he would not desert the African carriers who had accompanied him, right
across the continent to Quilimane on the east coast. On this journey he showed all the qualities of
a great explorer. His manner with the Africans was so patient that he never had to use violence. And
his scientific and geographical observations were minutely accurate. This one trip opened the heart
of Africa to the modern age.

But Livingstone was at all times more than a traveler. His cause was the gospel. His journal abounds
in passages of almost mystical devotion. Shortly before setting out on his great journey he wrote,
“I place no value on anything I have or may possess, except in relation to the kingdom of Christ.”

What moved him more than anything else was what he called “this open sore of the world,” the
devastating slave trade of central Africa. Speaking to the students at Cambridge in 1857 he said, “I
go back to Africa to try to make an open path for commerce and Christianity: do your best to carry
out the work which I have begun, I leave it to you.”

Commerce and Christianity? Was Livingstone simply a forerunner of those colonialist exploiters who
made life in so many parts of Africa a nightmare? No, far from it. Livingstone realized that the
slave trade could not continue apart from the African’s own participation in it. When slave raiding
was the way to wealth, the temptation was always present to engage in those raids on weaker
neighboring tribes that made life perilous for the defenseless. Only if the Africans could be
persuaded to engage in legitimate commerce, exchanging the products of their own fields and forests
for those desirable things the white man could supply, would the evil and destructive commerce be
brought to an end. That, at any rate, was Livingstone’s conviction, a central part of his dream for
Africa.

How was this missionary vision turned to action? What were the channels for this burst of spiritual
energy? The traditional denominations used one of three forms of church government: episcopal,
presbyterian, or congregational. The supporters of each claimed to be following the Bible, and all
the main arguments on each side were well-known. Men had suffered, some even shed blood, for each
form of government.

But as the conviction of the responsibility to spread the gospel worldwide began to dawn on British
and American Christians it became clear that none of the traditional forms of church government
would enable the church to embark on a world mission. Supporters of global evangelism were driven to
find another form of cooperation: the voluntary society.

\section{CREATION OF THE VOLUNTARY SOCIETY}

Again, Carey proved the pioneer. When he wrote his Enquiry he asked, what would a trading company
do? From this he proposed the formation of a company of serious-minded Christians, laymen and
ministers. The group should have a committee to collect and sift information and to find funds and
suitable missionaries to send to foreign lands.

The voluntary society, of which the missionary society was one early form, transformed
nineteenth-century Christianity. It was invented to meet a need rather than for theological reasons,
but in effect it undermined the established forms of church government. It made possible
interdenominational action. Anglican, Baptist, Congregationalist, and Methodist could work together
for defined purposes without raising troublesome questions of church structure. It also altered the
power base in the church by encouraging lay leadership. Ordinary Christian men, and later women,
came to hold key positions in the important societies, something thought impossible elsewhere in the
church.

These features appeared early in the history of missionary societies. The London Missionary Society
adopted in 1795 its “fundamental principle that our design is not to send Presbyterianism,
Independency, Episcopacy or any other form of church government~.~.~. but the glorious gospel of the
blessed God to the heathen.” One of the founders called for “the funeral of bigotry.”

Thanks to the creation of the societies, the missionary enthusiast who collected a penny a week from
members of his local missionary society auxiliary and distributed the missionary magazine
participated fully in the work of missions. Through the labors of such people, candidates for
service often came forward. The American missionary, Rufus Anderson, wrote in 1834, “It was not
until the present century that the evangelical churches of Christendom were ever really organized
with a view to the conversion of the world.” They became organized by means of the voluntary
society.

In the United States the first foreign missionary society was the American Board of Commissioners
for Foreign Missions (1810). It was formed on the initiative of a group of students at the newly
created Congregational Andover Theological Seminary. The leader of the group was Samuel J. Mills,
who while at Williams College had instigated the formation of a small society in which each member
pledged to devote his life to missionary service. At Andover, the society included the later famous
missionary to Burma, Adoniram Judson. Students, then, as so many times in later years, laid the
foundation for a missionary advance, this one for the youthful American nation.

In a few years Baptists, Presbyterians, and other major denominations followed the
Congregationalists in creating missionary agencies. The conversion of the heathen became one of the
major concerns of local congregations in every city and town in the country, stimulated by the
continuous activity of local societies and women’s organizations, “children’s days for foreign
missions,” occasional visits from missionaries on furlough, periodical campaigns for offerings, and,
more recently, the inclusion of support of foreign missions as a large item in regular church
budgets.

By the end of the nineteenth century, almost every Christian body, from the Orthodox Church of
Russia to the Salvation Army, and almost every country, from the Lutheran Church of Finland and the
Waldensian Church of Italy to the newest denomination in the United States, had its share in the
missionary enterprise overseas.

Many times these early missionaries were unaware of the conflicts that the gospel produced in other
cultures. To most of them, Christianity in its Western form was Christianity. Therefore to make an
Indian or a Malaysian a Christian was in great measure to turn him into a Dutchman or a Portuguese.

It is easy today to condemn this attitude. Yet every Christian society, and every individual
Christian, combines with the faith much that is cultural tradition. The problem has pursued all
missionary work from the beginning. Nevertheless, the consequences of such an alien presentation can
be disastrous to the progress of the gospel. If Christianity appears in Western habits that other
people find shocking, such as eating of meat or a greater familiarity between the sexes than is
permitted in most Eastern societies, then the faith is condemned before it is even examined.

On the other hand there is a distinctiveness of the Christian community that arises from the
distinctiveness of the gospel itself. The gospel is a revolutionary power and any attempt to
disguise this fact is likely to change the Christian faith into something else.

\section{THE MARKS OF MODERN CHRISTIANITY}

In spite of all barriers, however, the Protestant missionary movement continued to expand. In the
process it highlighted several characteristics of modern Christianity:

First, the worldwide expansion usually came by voluntaryism not by compulsion. Since the time of
Constantine the propagation of the faith had had the active support of rulers in Christian countries
and often resulted in mass conversions. Protestant missionaries in the nineteenth century, however,
went, with few exceptions, without state support or state control. They advanced only by the power
of persuasion. Thus Christians found a way to engage in missions while upholding religious freedom.

Second, this missionary movement tapped the wealth and talents of rank and file Protestants. Unlike
earlier missionary expansion led by monks and bishops, the new societies were organized on the
widest base possible. So the primary task of the churches, preaching the gospel, fell upon the
people of the churches.

Third, a wide variety of humanitarian ministries accompanied the widespread preaching of the gospel.
Mission agencies established schools, hospitals, and centers for training nurses and doctors. They
reduced many languages and dialects to writing and translated not only the Bible but other Western
writings into these languages. And they introduced public health measures and better agricultural
techniques. In some cases these activities were closely related to the goal of conversion, but many
sprang simply out of the recognition of social and physical needs that no Christian could in good
conscience ignore.

In many respects, then, the missionary movement restored the gospel to its central place in
Christianity. And in this important sense the movement recovered an element in the concept of the
holy catholic church that the splintering of the Reformation had obscured. A catholicity or
universality that inspired Carey’s workshop map reached out to embrace new peoples in many new
lands.

\readinghead{Suggestions for Further Reading}
\begin{readinglist}
\entry Drewery, Mary. William Carey: Shoemaker and Missionary. London: Hodder and Stoughton, 1978.
\entry George, Timothy. Faithful Witness: The Life and Mission of William Carey. Birmingham, AL: New
  Hope, 1991.
\entry Gascoigne, Bamber. The Christians. New York: William Morrow, 1977.
\entry Huxley, Elspeth. Livingstone and His African Journeys. New York: Saturday Review Press, 1974.
\entry Neill, Stephen. A History of Christian Missions. Edited by Owen Chadwick. Middlesex: Penguin
  Books, 1986.
\entry Northcott, Cecil. David Livingstone: His Triumph, Decline and Fall. Philadelphia:
  Westminster, 1973.
\end{readinglist}
